\section{The integer controlling the Frobenius trace differences}
\label{sec:common-trace-integer}

Fix an imaginary quadratic field \(K\), a positive integer \(N\),
a twist \(e\ge1\) and a prime \(p\ge7\) prime to
\(N\operatorname{disc}(K)\). Let \(w\mid p\), let
\(h=[\mathcal O_K/w:\mathbb F_p]\in\{1,2\}\), and set
\(q=p^h\) and \(Q=q^e\). Let \(O\) be a finite index set with
roots of unity \(z_o\in\mathbb C\) such that every field automorphism
of \(\mathbb C\) permutes the indexed values. Repetitions are
allowed and their multiplicities matter. Write
\[
                T_o=Q^{-1}z_o+1+Qz_o^{-1}.
\]
For the orbit in Section~\ref{sec:rh-plus-nonzero}, \(O\) consists
of the characters \(\xi\) with the prescribed kernel and
\(z_\xi=\xi(q)^{-1}\).

\subsection{The product and its size}

A finite family of trace values consists of a finite set \(S\)
and numbers \(t_s\in\mathbb C\) such that every automorphism of
\(\mathbb C\) permutes the indexed values and
\[
           |t_s|\le\sqrt q+2,\qquad
           Qt_s\ \text{is an algebraic integer}.
\]
At a coefficient prime \(\ell\nmid q\) and for a field
isomorphism \(\iota:\overline{\mathbb Q}_\ell\to\mathbb C\), a
family of admissible representations is \emph{classified} by \(S\)
at \(g\) if each \(\rho\) has a label \(s(\rho)\in S\) with
\[
 \iota\bigl(\operatorname{tr}\rho(g)\bigr)=t_{s(\rho)},\qquad
 s(\rho)=s(\rho')\ \Longrightarrow\ \rho\simeq\rho'.
\]
Several labels can have equal numeric traces. The implication
about labels preserves the multiplicity of their contributions.

Define the product before choosing any coefficient prime,
representations or isomorphism by
\begin{equation}
   n_p=\prod_{o\in O}\prod_{s\in S} F_{o,s},
   \qquad
   F_{o,s}=Q(t_s-T_o)=Qt_s-z_o-Q-Q^2z_o^{-1}.
   \label{eq:common-trace-product}
\end{equation}
We claim that \(n_p\in\mathbb Z\setminus\{0\}\), that
\begin{equation}
  \log|n_p|\le |O|\,|S|\,(2e\log q+1),
  \label{eq:common-trace-log}
\end{equation}
and that for every classified family and every finite selection
\(X_o\) of pairwise nonisomorphic members for each \(o\),
\begin{equation}
 |n_p|_\ell\le
 \prod_{o\in O}\prod_{\rho\in X_o}
 \left|\operatorname{tr}\rho(g)-\iota^{-1}(T_o)\right|_\ell .
 \label{eq:common-trace-subproduct}
\end{equation}

To separate \(t_s\) and \(T_o\), note that \(Q\ge q\ge7\) and
\(\sqrt q\le(q+9)/6\) by \((\sqrt q-3)^2\ge0\).
It follows that \(Q-\sqrt q-3\ge4/3>Q^{-1}\). By the reverse
triangle inequality,
\[
              |T_o|\ge Q-1-Q^{-1}>\sqrt q+2 .
\]
Every factor of \eqref{eq:common-trace-product} is therefore
nonzero. Each factor is integral because \(Qt_s\), \(z_o\),
\(z_o^{-1}\) and \(Q\) are algebraic integers. In addition,
\[
 |F_{o,s}|\le Q^2+Q(\sqrt q+3)+1\le2Q^2 .
\]
The last inequality follows from \(Q(Q-\sqrt q-3)>1\).

All factors belong to a finite normal extension of \(\mathbb Q\)
in \(\mathbb C\). Each automorphism of this extension extends to an
automorphism of \(\mathbb C\) by extending first to an algebraic
closure and then using a transcendence basis. The resulting
permutations of \(O\) and \(S\) jointly permute \(O\times S\).
Their product \(n_p\) is consequently rational. A rational algebraic
integer is an integer. The empty product is \(1\). Taking logarithms
and using \(\log2\le1\) proves \eqref{eq:common-trace-log}.

Finally transport all the factors by \(\iota^{-1}\).
Every resulting algebraic integer has absolute value at most one
in \(\overline{\mathbb Q}_\ell\). For each \(o\), the map
\(\rho\mapsto s(\rho)\) is injective on \(X_o\). Discarding the
factors whose labels are outside its image can only increase
the \(\ell\)-adic absolute value of the product. Since
\(|Q|_\ell=1\), the remaining factors give
\eqref{eq:common-trace-subproduct}. In particular every trace
difference in its right side is nonzero. This proves all three
assertions for arbitrary trace data with the stated hypotheses.

\subsection{Trace data from the two kinds of eigensystems}

We describe exactly which local estimates give these hypotheses.
The rank-three CM eigensystems for the inverse cyclotomic twist
have three parameters \(b_i\) with the convention that the
eigenvalues of the twisted representation at \(\sigma^{-1}\) are
the \(b_i^{-1}\), where \(\sigma\) is arithmetic Frobenius at \(w\).
The weight is \(-2\) and the integral weight vector is
\((e-1,-1,-e-1)\). The local estimates on the parameters are
\[
     q^{-1}q^{-1/2}< |b_i| < q^{-1}q^{1/2},\qquad
     \prod_{i=1}^3 |b_i|=q^{-3},\qquad
     q^{e-1}b_i^{-1}\ \text{is integral}.
\]
Here integrality means algebraic integrality over \(\mathbb Z\).
Undoing the twist divides each eigenvalue by \(q\), since
\(\chi_\ell(\sigma)=q\). The trace label for the original
three-dimensional representation is
\[
                         t_s=\sum_{i=1}^3(qb_i)^{-1}.
\]
Thus \(Qt_s\) is integral. Put \(a=\sqrt q\) and
\(x_i=|(qb_i)^{-1}|\). We have \(1/a\le x_i\le a\) and
\(\prod_i x_i=1\). Reorder them so that \(x_1\ge x_2\ge x_3\).
If \(x_2\le1\), their sum is at most \(a+2\).
If \(x_2\ge1\), then \(x_1+x_2\le x_1x_2+1\) and
\(x_1x_2=x_3^{-1}\le a\), while \(x_3\le1\).
The same bound holds in both cases, so \(|t_s|\le\sqrt q+2\).

For the rank-two systems over \(\mathbb Q\), the prime has residue
size \(p\), the weight is zero and the integral weight vector is
\((e,-e)\). Their two parameters satisfy
\[
   p^{-1/2}<|b_i|<p^{1/2},\qquad |b_1b_2|=1,\qquad
             p^e b_i^{-1}\ \text{is integral}.
\]
The representation over \(K\) is isomorphic to
\(\mathbf1\oplus W|_{G_K}\). On the unramified quotient at the
chosen place, arithmetic Frobenius over \(K\) restricts to the
\(h\)-th power of arithmetic Frobenius over \(\mathbb Q\).
Consequently its inverse gives
\[
                   t_s=1+b_1^{-h}+b_2^{-h}.
\]
The identity \(Qb_i^{-h}=(p^eb_i^{-1})^h\) gives the claimed
integrality. Of the two numbers \(|b_1|^{-h}, |b_2|^{-h}\) with
product one, the larger is at most \(\sqrt q\) and the smaller
at most one. This again gives \(|t_s|\le\sqrt q+2\).
In both ranks the operations used to make \(t_s\) commute with
field automorphisms. A permutation of the parameter multisets
therefore permutes the trace labels.

The absolute value bounds here are precisely the Satake bounds
used in the released argument. For a unitary cuspidal
representation of \(\mathrm{GL}_n\) at a spherical place of residue
size \(u\), the Jacquet--Shalika estimate puts the absolute values
of its normalized Satake parameters strictly between \(u^{-1/2}\)
and \(u^{1/2}\). Their product has absolute value one by unitarity
of the central character. Twisting to weight zero gives the
displayed inequalities after reinstating the normalizing powers of
\(q\) or \(p\). See H.~Jacquet and J.~Shalika, \emph{On Euler
products and the classification of automorphic representations I},
Amer.~J.~Math.~103 (1981), 499--558. Integrality of the scaled
inverse parameters is the separate integral Satake estimate in the
source.

\subsection{Global labels and the quadratic restriction argument}

A label for an automorphic representation over a CM field is its
full system of normalized unramified Hecke multisets away from
the prescribed level. The test place \(w\) is only one entry.
Suppose the relevant cuspidal representations at weight vector
\(m\) and level \(M\) have at most \(C(M(H+1))^A\) pairwise
nonisomorphic members, where \(H\ge\max_i|m_i|\). Suppose also
that the system depends only on the isomorphism class, its Galois
conjugates occur in the same collection, and the attached
semisimple Galois representations agree with it at places
away from \(M\ell\). The set of distinct global systems then has
the same upper bound. Indeed any finite subset of the image lifts
to pairwise nonisomorphic preimages by choosing one form for each
system. Thus it is a finite set of labels independent of \(\ell\).
Conjugation permutes these labels. Equality of labels gives equality
of the characteristic polynomials of the attached representations
at all but finitely many Frobenius classes. By Chebotarev these
classes are dense up to conjugacy. Continuity gives equality of
traces everywhere, and the characteristic-zero Brauer--Nesbitt
theorem gives isomorphism for semisimple representations.
These statements construct exactly the classification used above
when the indicated automorphic estimates and compatibility hold.

For the rational rank-two systems the source first constructs
finite collections over auxiliary imaginary quadratic fields. Here
is the finite counting argument allowing those collections to
classify the rational systems. Choose six imaginary quadratic
fields \(K_i\) of pairwise coprime discriminants. For example
one may take \(K_i=\mathbb Q(\sqrt{-d_i})\) with
\(d_i=3,7,11,19,23,31\). At most one of these discriminants is
divisible by the test prime \(p\), leaving five choices. Write
\(\delta_i\) for the distinct quadratic characters of
\(G_{\mathbb Q}\) with kernels \(G_{K_i}\).

For an irreducible two-dimensional representation \(W\) over
\(\overline{\mathbb Q}_\ell\), at most four distinct characters
\(\delta\) satisfy \(W\simeq W\otimes\delta\). Each such
intertwiner is a nonzero vector in the four-dimensional
\(\operatorname{End}(W)\) on which conjugation acts by
\(\delta\) (or by its inverse, equal to it here). Nonzero
eigenvectors belonging to distinct characters of a group are
linearly independent. Indeed a minimal relation with at least
two terms yields a shorter nonempty relation by subtracting a
scalar multiple of its translate by an element distinguishing
two of the characters. Consequently one of the five remaining
\(i\) satisfies \(W\not\simeq W\otimes\delta_i\).

For \(H_i=G_{K_i}\), restriction to \(H_i\) is semisimple on
this irreducible \(W\). In fact choose a simple \(H_i\)-submodule
and sum its two conjugates under coset representatives. The span
is nonzero and \(G_{\mathbb Q}\)-stable, hence all of \(W\), and
a sum of simple submodules is semisimple. Frobenius reciprocity
and \(\operatorname{Ind}_{H_i}^{G_{\mathbb Q}}\mathbf1
\simeq\mathbf1\oplus\delta_i\) give
\[
 \operatorname{End}_{H_i}(W)
    \simeq\operatorname{End}_{G_{\mathbb Q}}(W)
       \oplus\operatorname{Hom}_{G_{\mathbb Q}}
                             (W,W\otimes\delta_i)
\]
as vector spaces. Its dimension is one for our chosen \(i\), so
\(W|_{H_i}\) is irreducible. The same adjunction for two irreducible
representations \(W,W'\) shows that isomorphic restrictions force
\(W'\simeq W\) or \(W'\simeq W\otimes\delta_i\). Hence at most two
isomorphism classes of rational representations occur over a fixed
restricted class. The arithmetic hypothesis presently needed is
that an admissible rational system with irreducible restriction to
\(H_i\) belongs there to the indicated class of polarized CM systems.

Choose a place above \(p\) in each of these five fields and let its
residue degree be \(f_i\), equal to one or two. From a restricted
label \(s\in S_i\) and its parameter multiset \(B_i(s)\), form all
two-element multisets \(B\) such that
\(\{b^{f_i}:b\in B\}=B_i(s)\), with multiplicity. For each \(s\)
there are at most \(f_i^2\le4\) possibilities. Enumerate the two
target entries, choose one of their \(f_i\)-th roots at each
position, and forget the order. Append two copies of each label
\((i,s,B)\). The number of labels is at most
\[
                       8\sum_i |S_i|.
\]
Galois conjugation permutes these labels, acting on \(s\) and \(B\)
and leaving the copy index fixed. The bounds on the absolute values
and their product pass to \(f_i\)-th roots. The needed integrality
passes as well, since \((p^e b^{-1})^{f_i}\) integral implies
\(p^eb^{-1}\) integral by transitivity of integrality.

To classify a rational representation, choose one of its irreducible
restrictions as above, take the corresponding restricted label, and
choose the actual parameter multiset at \(p\). The power relation at
Frobenius ensures that it is among the listed lifts. At each fixed
triple \((i,s,B)\) assign distinct copy indices to its at most two
underlying rational isomorphism classes. This gives classification
for each \(\ell\) and \(\iota\). The two-copy label set used by all
\(\ell\) and \(\iota\) was fixed before making these assignments.
The automorphic restriction assertion, the estimates for the
polarized CM class, and the level and newform criteria which supply
the requisite rank-two representation are inputs to this reduction.

\subsection{The polynomial bound uniform in the prime of coefficients}

Write \(D_K=|\operatorname{disc}K|\) and \(Y=N(e+1)\).
Suppose the two global classifications above yield label sets
\(S_1,S_2\), with
\[
             |S_1|\le c_1Y^{A_1},\qquad
             |S_2|\le c_2(D_KY)^{A_2}.
\]
For \(S_1\) one uses level \(N\). The weight bound \(H=e+1\)
would give \(N(e+2)\), and \(e+2\le2(e+1)\) allows the extra
factor \(2^{A_1}\) to be absorbed into \(c_1\).
For \(S_2\) the rational level is \(ND_K\) and \(H=e\).
The sets fit together by disjoint union \(S=S_1\sqcup S_2\).
If a representation satisfies both descriptions, give it its
first label. Equality of disjoint-union labels then still implies
isomorphism. We obtain
\[
 |S|\le(c_1+c_2D_K^{A_2})Y^{A_1+A_2},\quad
 2e\log q+1\le 5Y\log p ,
\]
using \(q=p^h,\ h\le2,\ e+1\le Y\) and \(1<\log p\).
Equations~\eqref{eq:common-trace-log} and
\eqref{eq:common-trace-subproduct}, with \(g=\sigma^{-1}\),
give the required common integer with
\[
 \log|n_p|\le 5(c_1+c_2D_K^{A_2})\,
                 |O|\,Y^{A_1+A_2+1}\log p .
\]
Both the size bound and the simultaneous subproduct inequality
apply on the same \(e,N,p,w,O\). Only the subcollections of
representations vary with the coefficient prime.

\paragraph{Dependencies of the assertion for the actual families.}
The release's \texttt{number\allowbreak{}U3\allowbreak{}U\_\allowbreak{}of\_\allowbreak{}prints} uses the CM
and rational counting arguments in the dense-threshold file
\texttt{Frobenius\allowbreak{}Trace\allowbreak{}Count\allowbreak{}Bound\allowbreak{}Of\allowbreak{}Cm\allowbreak{}Eigensystem\allowbreak{}Count}.
Its inputs include the polynomial count of regular algebraic
cuspidal systems of the given level and weights, stability under
field automorphisms, the integral Satake bounds, and compatibility
of the admitted Galois representations with the global Hecke
systems. On the rational side it also uses the polarized CM
class counts, passage to irreducible quadratic restrictions, and
the newform-admissibility comparison. The local parameter estimates
for these two families and the common-product theorem occur in
\texttt{Admissible\allowbreak{}Residues.\allowbreak{}Frobenius\allowbreak{}Trace\allowbreak{}Count\allowbreak{}Bound\allowbreak{}Of\allowbreak{}Rat\allowbreak{}Eigensystem\allowbreak{}Count}.
The archimedean bound in the associated trace data is
\(\sqrt q+2\). The preceding construction proves the common integer
from the stated automorphic inputs. The separate deformation
inequality for one Selmer character, denoted by
\texttt{One\allowbreak{}Char\allowbreak{}U\_\allowbreak{}s}, is used in Section~\ref{sec:rh-plus-nonzero}.
